\chapter{Preface}

《数据库系统实现》是电子科技大学为本科同学开设的一门挑战课，内容主要介绍数据库系统的内部实现。

从数据库课程体系来看，数据库系统的课程主要分为几个级别：

\begin{itemize}
    \item 数据库系统原理：介绍数据库系统的基本概念、数据模型与应用系统设计。
    \item 数据库系统实现：详细介绍数据库系统的内部实现，包括存储管理、查询处理、事务管理等。
    \item 数据库系统性能优化：讨论数据库系统的性能优化技术，如索引优化、查询优化等。
    \item 分布式数据库系统：介绍分布式数据库系统的基本概念、原理和实现。
\end{itemize}

目前，国内本科阶段的数据库系统教学主要集中在原理方面的讲述，开设数据库系统实现的课程比较少。本课程受学校资助，面向本科同学介绍数据库系统的内部实现，内容包括存储管理、查询处理、事务管理等。

课程开设的时间比较早，至今已有七八个年头。课程一直采用Stanford教材\cite{garcia2000database}，这是国内能够找到的为数不多的详细介绍数据库内核的教材。但在教学中发现，教材虽经更新，内容上仍然与工程上的实际情况有一定差距，例如存储管理部分对底层文件、记录的介绍不够详细等。有鉴于此，撰写了这本讲义，重新组织各章内容。

